% GENERATED FILE. Edit canonical lessons, then run npm run generate:books.
% canonical-source-hash: fnv1a64:07a84c3fa285eb90
% canonical-lessons: LA-C91-avis, LA-C91-lupus, LA-C91-leo, LA-C91-ursus, LA-C91-serpens

\chapter{Bird, Wolf, Lion, Bear, Snake}
\label{ch:la-bird-wolf-lion-bear-snake}
\hlchaptermodality{91}

\begin{chapteropening}
\textbf{By the end of this chapter:} \emph{I can name a bird, a wolf, a lion, a bear and a snake.}

Complete the last lesson of chapter 91: I can name a bird, a wolf, a lion, a bear and a snake.
\end{chapteropening}

\section[avis, avis]{avis, avis — a bird}

\label{lesson:LA-C91-avis}



\begin{quote}
Before the new one: say the Latin for a goat, then the Latin for a pig.
\end{quote}

\subsection*{You'll want to know: avis, avis}
\textbf{avis, avis} — \textquotedblleft{}a bird\textquotedblright{}. English \emph{aviation}.

\begin{cousinweb}
\textbf{avis} is inherited, usually traced to Proto-Indo-European \emph{*h2ewis}, \textquotedblleft{}bird\textquotedblright{} (EtymDB 2.1). \emph{auspex}, \textquotedblleft{}a bird-watcher\textquotedblright{}, the priest who read omens from birds, is built on it.
\end{cousinweb}

\begin{grammarlens}[title={What you've built}]
One more word for this chapter.
\end{grammarlens}

\subsection*{Your turn}
\begin{itemize}
\raggedright
  \item \emph{Say it:} \emph{avis}
  \item \emph{Say it:} \emph{avis}, once more
  \item \emph{Say it:} say \emph{avis}
  \item \emph{Recall:} say the Latin for a goat, then the Latin for a pig, then say \emph{avis} again
\end{itemize}

\begin{tcolorbox}[breakable,colback=teal!4,colframe=teal!35!black,title={Before you move on}]
What does avis, avis mean? (\textquotedblleft{}A bird\textquotedblright{}.) Say it once more.
\end{tcolorbox}
\section[lupus, lupī]{lupus, lupī — a wolf}

\label{lesson:LA-C91-lupus}



\begin{quote}
Before the new one: say the Latin for a pig, then the Latin for a bird.
\end{quote}

\subsection*{You'll want to know: lupus, lupī}
\textbf{lupus, lupī} — \textquotedblleft{}a wolf\textquotedblright{}.

\begin{cousinweb}
\textbf{lupus} is usually traced to Proto-Indo-European \emph{*wlkwos}, \textquotedblleft{}wolf\textquotedblright{} (EtymDB 2.1), with Greek \emph{lykos}, \textquotedblleft{}wolf\textquotedblright{}, and English \emph{wolf} as cousins (Lewis \& Short). Its p, where Latin would have qu, suggests it came in from a neighbouring Italic language such as Sabine.
\end{cousinweb}

\begin{grammarlens}[title={What you've built}]
One more word for this chapter.
\end{grammarlens}

\subsection*{Your turn}
\begin{itemize}
\raggedright
  \item \emph{Say it:} \emph{lupus}
  \item \emph{Say it:} \emph{lupus}, once more
  \item \emph{Say it:} say \emph{lupus}
  \item \emph{Recall:} say the Latin for a pig, then the Latin for a bird, then say \emph{lupus} again
\end{itemize}

\begin{tcolorbox}[breakable,colback=teal!4,colframe=teal!35!black,title={Before you move on}]
What does lupus, lupī mean? (\textquotedblleft{}A wolf\textquotedblright{}.) Say it once more.
\end{tcolorbox}
\section[leō, leōnis]{leō, leōnis — a lion}

\label{lesson:LA-C91-leo}



\begin{quote}
Before the new one: say the Latin for a bird, then the Latin for a wolf.
\end{quote}

\subsection*{You'll want to know: leō, leōnis}
\textbf{leō, leōnis} — \textquotedblleft{}a lion\textquotedblright{}.

\begin{cousinweb}
\textbf{leō} was borrowed from Greek \emph{leōn}, \textquotedblleft{}lion\textquotedblright{} (Lewis \& Short); where Greek got the word is uncertain. English \emph{lion} comes from it through French.
\end{cousinweb}

\begin{grammarlens}[title={What you've built}]
One more word for this chapter.
\end{grammarlens}

\subsection*{Your turn}
\begin{itemize}
\raggedright
  \item \emph{Say it:} \emph{leō}
  \item \emph{Say it:} \emph{leō}, once more
  \item \emph{Say it:} say \emph{leō}
  \item \emph{Recall:} say the Latin for a bird, then the Latin for a wolf, then say \emph{leō} again
\end{itemize}

\begin{tcolorbox}[breakable,colback=teal!4,colframe=teal!35!black,title={Before you move on}]
What does leō, leōnis mean? (\textquotedblleft{}A lion\textquotedblright{}.) Say it once more.
\end{tcolorbox}
\section[ursus, ursī]{ursus, ursī — a bear}

\label{lesson:LA-C91-ursus}



\begin{quote}
Before the new one: say the Latin for a wolf, then the Latin for a lion.
\end{quote}

\subsection*{You'll want to know: ursus, ursī}
\textbf{ursus, ursī} — \textquotedblleft{}a bear\textquotedblright{}.

\begin{cousinweb}
\textbf{ursus} is inherited, usually traced to Proto-Indo-European \emph{*h2rtkos}, \textquotedblleft{}bear\textquotedblright{} (EtymDB 2.1); Greek \emph{arktos}, \textquotedblleft{}bear\textquotedblright{}, is a cousin (Lewis \& Short). English \emph{Arctic}, the region under the stars of the Bear, comes from the Greek.
\end{cousinweb}

\begin{grammarlens}[title={What you've built}]
One more word for this chapter.
\end{grammarlens}

\subsection*{Your turn}
\begin{itemize}
\raggedright
  \item \emph{Say it:} \emph{ursus}
  \item \emph{Say it:} \emph{ursus}, once more
  \item \emph{Say it:} say \emph{ursus}
  \item \emph{Recall:} say the Latin for a wolf, then the Latin for a lion, then say \emph{ursus} again
\end{itemize}

\begin{tcolorbox}[breakable,colback=teal!4,colframe=teal!35!black,title={Before you move on}]
What does ursus, ursī mean? (\textquotedblleft{}A bear\textquotedblright{}.) Say it once more.
\end{tcolorbox}
\section[serpēns, serpentis]{serpēns, serpentis — a snake}

\label{lesson:LA-C91-serpens}



\begin{quote}
Before the new one: say the Latin for a lion, then the Latin for a bear.
\end{quote}

\subsection*{You'll want to know: serpēns, serpentis}
\textbf{serpēns, serpentis} — \textquotedblleft{}a snake\textquotedblright{}. Literally \textquotedblleft{}the creeper\textquotedblright{}.

\begin{cousinweb}
\textbf{serpēns} is the present participle of \emph{serpō}, \textquotedblleft{}to creep\textquotedblright{} (Lewis \& Short), used as a noun: the creeping thing. \emph{serpō} is usually traced to Proto-Indo-European \emph{*serp-}, \textquotedblleft{}to crawl, to creep\textquotedblright{} (EtymDB 2.1).
\end{cousinweb}

\begin{grammarlens}[title={What you've built}]
One more word for this chapter.
\end{grammarlens}

\subsection*{Your turn}
\begin{itemize}
\raggedright
  \item \emph{Say it:} \emph{serpēns}
  \item \emph{Say it:} \emph{serpēns}, once more
  \item \emph{Say it:} say \emph{serpēns}
  \item \emph{Recall:} say the Latin for a lion, then the Latin for a bear, then say \emph{serpēns} again
\end{itemize}

\begin{tcolorbox}[breakable,colback=teal!4,colframe=teal!35!black,title={Before you move on}]
What does serpēns, serpentis mean? (\textquotedblleft{}A snake\textquotedblright{}.) Say it once more.
\end{tcolorbox}
